\documentclass[10pt,a4paper]{article}
\usepackage[left=2.5cm,right=2.5cm,top=2.5cm,bottom=2.5cm,]{geometry}
\usepackage{amsmath,amssymb,amsthm}
\usepackage{mathrsfs}
\newtheorem{theorem}{Theorem}[section]
\newtheorem{lemma}[theorem]{Lemma}
\newtheorem{proposition}[theorem]{Proposition}
\usepackage{setspace}
\setlength{\parindent}{0pt}
\onehalfspacing
%chord
\newcommand{\sh}{^\sharp}
\newcommand{\fl}{\flat}
\begin{document}
\subsubsection{12-tone-equal temperament}
In physics, the essence of sound is a mechanical wave. Taking the 12-tone equal temperament as an exapmle,
In this system, the ratio between adjacent semitones is constant $2^{\frac{n}{12}}$, calculate equation $f=2^{\frac{n}{12}}f_0$, 
$\mathscr{P}=\{A,A^{\sharp},B_{\flat},B,C,C^{\sharp},D_{\flat},D,D^{\sharp},E_{\flat},E,F,F^{\sharp},G_{\flat},G,G^{\sharp},A_{\flat},A\}$,
we can use $D_n$ to describe.
\end{document}



